\subsection{Recuperación de contraseña con activación explícita}\label{sec:impl-recuperacion-password}

La recuperación realiza el flujo alterno de CU-03 y la parte correspondiente de RF-01 (\cref{tab:cu-03,tab:rf-acceso}): un enlace temporal por correo permite fijar otra contraseña (\cref{fig:impl-recuperacion-estados}). Es distinta del código de recuperación MFA, que exige conocer la contraseña, y sus rutas, formulario y correo están deshabilitados por omisión (ADR 0066).

\begin{figure}[!htb]
  \centering
  \begin{tikzpicture}[
  est/.style={cajaClara, text width=2.5cm, minimum height=1.15cm, rounded corners=5pt},
  ev/.style={etq, align=center}]

  % Estados principales
  \node[est] (s0) at (1.4,-1.9) {\textbf{Solicitada}\\ respuesta 202; no espera búsqueda, emisión ni envío};
  \node[est] (s1) at (5.65,-1.9) {\textbf{Emitida}\\ PostgreSQL guarda solo el SHA-256; el correo lleva el enlace};
  \node[est] (s2) at (9.9,-1.9) {\textbf{Consumida}\\ contraseña nueva; la capacidad ya no sirve};
  \node[cajaOsc, text width=2.9cm, minimum height=1.15cm] (fin) at (14.35,-1.9)
    {Acceso con la contraseña nueva y el MFA ya enrolado};

  \draw[flecha] (s0) -- node[ev, above=3pt]{cuenta\\ activa y\\ cupo libre} (s1);
  \draw[flecha] (s1) -- node[ev, above=3pt]{capacidad\\ vigente y\\ contraseña} (s2);
  \draw[flecha] (s2) -- node[ev, above=3pt]{sin\\ sesión} (fin);

  % Entrega incierta
  \node[paso, text width=3.4cm, anchor=south] (inc) at (5.65,-0.95)
    {Entrega incierta: la capacidad sigue vigente, sin reenvío automático};
  \draw[flechaSuave] (inc.south) -- (s1.north);

  % Salidas
  \node[cajaGris, text width=2.5cm, minimum height=1.0cm] (n0) at (1.4,-3.95)
    {\textbf{Sin emisión}: cuenta desconocida o inactiva, o cupo ocupado; mismo resultado, sin otra tarea};
  \node[cajaError, text width=1.9cm, minimum height=1.0cm] (e1) at (4.35,-3.95) {\textbf{Vencida}\\ (vigencia\\ configurada)};
  \node[cajaError, text width=2.3cm, minimum height=1.0cm] (e2) at (6.85,-3.95) {\textbf{Cancelada}: entrega rechazada (mismo ID y digest) o restauración};
  \draw[flecha] (s0) -- (n0);
  \draw[flechaError] (s1.south -| e1) -- (e1);
  \draw[flechaError] (s1.south -| e2) -- (e2);

  % Instalación privada
  \node[paso, text width=6.6cm, anchor=north west] at (9.0,0.05)
    {\textbf{Instalación privada}: Resend; vigencia de 900 s; un enlace pendiente por cuenta; una solicitud por correo cada 300 s; espera total de correo de 5000 ms};

  % Transacción de consumo
  \node[pasoOsc, text width=7.0cm, anchor=north] (tx) at (12.25,-2.85)
    {\textbf{Una transacción}, tras los bloqueos de auditoría y de fila:\\
     1. vuelve a comprobar vigencia y generación (la inspección previa a Argon2id no autoriza);\\
     2. sustituye el hash Argon2id de la contraseña;\\
     3. revisión +1 y generación de autenticación +1;\\
     4. consume la capacidad y agrega el evento.\\
     Si falla la auditoría, se revierte todo. Rol, pertenencias y MFA se conservan.};
  \draw[flechaSuave] (s2.south) -- (s2.south |- tx.north);
\end{tikzpicture}

  \caption{Ciclo de una capacidad de recuperación. La respuesta pública es igual haya o no emisión; el consumo avanza la generación (\cref{fig:impl-identidad-estado}).}
  \label{fig:impl-recuperacion-estados}
\end{figure}

\paragraph{Capacidad y transporte.} La capacidad es un valor aleatorio de 32 bytes del que se guarda solo una huella, con UUID propio por emisión. Un enlace nuevo no reemplaza otro vigente. La solicitud admite un correo acotado, y un único cupo cubre la entrada pendiente y su ejecución. El consumidor comparte con HTTP el presupuesto bloqueante, cancelar la espera no libera el permiso y, al detener el servidor, se descartan las tareas no iniciadas y se esperan las ya comenzadas. El emisor usa un destino HTTPS fijo de Resend, sin redirecciones, proxies de entorno ni reintentos, con la capacidad en el fragmento de una URL de origen aprobado.

\paragraph{Formulario.} Qadra retira el fragmento, pide confirmar la contraseña y no guarda el enlace; ante una respuesta incierta no repite el envío: se comprueba el acceso o se pide otro enlace, que descarta las respuestas tardías. Con sesión activa conserva los borradores (\cref{sec:impl-qadra-borradores}) antes de exigir otro acceso con MFA.

\paragraph{Configuración.} La activación exige habilitación expresa, remitente, origen, vigencia y cupo sin valor por omisión, cuotas, ventanas y tiempos del transporte; una clave de correo compartida no activa ningún canal, y los valores incompletos impiden abrir el servidor.

\paragraph{Límites.} Admitir una solicitud no es entrega durable ni prueba que la cuenta exista, y el resultado igual para cuentas desconocidas o inactivas no acredita uniformidad temporal del transporte. El rol operativo solo invoca operaciones restringidas, sin leer capacidades ni hashes. La habilitación operativa depende del transporte, de los límites y URL, de la retención y de la restauración: tras restaurar un respaldo, un comando administrativo (\cref{tab:impl-cli-catalogo}) que verifica la base, el esquema y la cabeza de auditoría esperados cancela las capacidades pendientes con un recibo, sin tocar credenciales, MFA ni capacidades consumidas. La verificación y la aceptación productiva están en la \cref{sec:pruebas-recuperacion-password}.

\Fuentes{\rutarepo{docs/password-reset-public.md}; ADR 0066.}
